\documentclass[10pt,a4paper]{amsart}
\usepackage[margin=19mm]{geometry}
\usepackage{amsmath,amssymb,mathtools}
\usepackage[scr=rsfs,cal=euler]{mathalfa}
\usepackage{xfrac}
\usepackage[unicode,hidelinks,bookmarks=false]{hyperref}
\newcommand{\ie}{i.e.\ }
\newcommand{\cf}{cf.\ }
\newcommand{\resp}{respectively\ }
\newcommand{\Subsection}{\subsection}
\providecommand{\nobreakdash}{\nobreak\hbox{-}}
\setlength{\emergencystretch}{2em}
\multlinegap0pt
\usepackage{bm}
\usepackage{bbm}
\usepackage{dsfont}
\usepackage{xspace}
\usepackage{cancel}
\usepackage{mathtools}
\usepackage{amsthm}
\makeatletter
\@ifundefined{theorem}{\newtheorem{theorem}{Theorem}}{}
\@ifundefined{lemma}{\newtheorem{lemma}[theorem]{Lemma}}{}
\makeatother
\title[The convergence of inversive distance circle packings to the]{The convergence of inversive distance circle packings to the Riemann mapping}
\author{Yuxiang Chen, Yanwen Luo, Xu Xu}
\date{Source: arXiv:2204.08145v2 (CC BY 4.0)}
\begin{document}
\maketitle

\noindent\emph{Source and licence.} The convergence of inversive distance circle packings to the Riemann mapping. Version arXiv:2204.08145v2 (CC BY 4.0), distributed under the Creative Commons Attribution 4.0 International licence, as recorded in the arXiv licence metadata for this exact version and shown on the abstract page https://arxiv.org/abs/2204.08145v2. Full rights notice: licenses/arxiv-2204.08145-CC-BY-4.0.txt.\par\smallskip

\noindent\textit{Editorial scope.} Counted unit: the Lemma whose source label is \texttt{ratio} in the article (source0.tex lines 888--897, source label \texttt{ratio}), delivered together with the article's own complete original proof of it (source0.tex lines 899--951, one \texttt{proof} environment of 53 lines). No mathematics is written, repaired or re-derived: statement and proof are the article's own text line for line. Required context retained uncounted: h\_ij,k (lines 511--517); h\_i (lines 520--525). Every definition, display and label of those lines is retained unchanged. Omitted from this package and not redistributed: 1--510, 518--519, 526--887, 898--898, 952--1406. Nothing inside a retained excerpt is omitted or altered: every retained body is delivered exactly as the article's lines are written, so no omission receipt of any kind accompanies this package. Citations: the retained excerpts carry no citation command at all, so no bibliography entry of the article is transcribed and no reference is redistributed. Reference adaptation: none was needed and none was made; every reference inside the delivered unit resolves inside it, so no unresolved reference or citation is produced. Printed slips kept literal: no printed word, symbol, hypothesis or number of the counted statement or of its proof was altered. Layout and class: the article's own class is \texttt{amsart} with packages including enumitem, psfrag, graphicx, pinlabel, color, amsmath, amssymb, amscd, picinpar, times, pb-diagram, wrapfig, xspace, hyperref; the wrapper is the lane's pinned \texttt{amsart} editorial wrapper at 10pt on A4, which restates the theorem-like environments the retained text needs and adds the article's own macro definitions that the retained text needs (none), with extra wrapper packages (bm, bbm, dsfont, xspace, cancel, mathtools). The delivered numbering is the unit's own. Assets: no retained span contains a figure environment, a graphics-inclusion command, a PDF-page inclusion, a table or a TikZ picture; no image file is copied into this package, the package declares no asset dependency and no image file is redistributed with it. Licence: version arXiv:2204.08145v2 of this article is distributed under the Creative Commons Attribution 4.0 International licence, as recorded in the arXiv licence metadata for this exact version and shown on the abstract page https://arxiv.org/abs/2204.08145v2. A full rights notice is delivered at licenses/arxiv-2204.08145-CC-BY-4.0.txt. Characters: every retained excerpt is pure ASCII, so the delivered engine, XeLaTeX with its default Latin Modern text font, reports no missing character.
		\begin{equation}\label{h_ij,k}
			\begin{aligned}
				h_{ij,k}
			=\frac{r_1^2r_2^2r_3^2}{A_{123}l_{ij}}[\kappa_k^2(1-I_k^2)+\kappa_j\kappa_k\gamma_{i}+\kappa_i\kappa_k\gamma_{j}]
				=\frac{r_1^2r_2^2r_3^2}{A_{123}l_{ij}}\kappa_kh_k
			\end{aligned}
		\end{equation}

		\begin{equation}
		\label{h_i}
			\begin{aligned}
h_k=\kappa_k(1-I_k^2)+\kappa_i\gamma_{j}+\kappa_j\gamma_{i}.
			\end{aligned}
		\end{equation}

\begin{lemma}
\label{ratio}
	Let $\triangle v_1v_2v_3$ be a weighted triangle generated by an inversive distance  circle packing $(r_1, r_2, r_3)$ and the weight $I>1$ is a constant. There exists a constant $\theta_0 = \theta_0(I)\in (0, \frac{\pi}{6})$ such that if the three inner angles of the triangle are bounded in $[\pi /6 - \theta_0,  \pi/2 + \theta_0]$, then
	\begin{enumerate}
		\item[(a)] $r_j/r_i \leq 20$ for any two radii $r_i$ and $r_j$,
		\item[(b)] $C\leq \eta_{ij}^k \leq M$ for some constants $C = C(I)>0$  and $M = M(I)>0$.
%		\item If the angles of the triangle $\triangle ijl$ satisfying the same condition, then
%		$\eta_{ij}^k/\eta_{ij}^l \geq R$ for some constant $R = R(I)>0$.
	\end{enumerate}
\end{lemma}

\begin{proof}
    Set $$\theta_{0} = \min \{\frac{\pi}{1000}, \arcsin \frac{1}{10(20+I)}\}.$$


	To prove part (a), by the angle bound and the sine law,
\begin{equation}\label{lij bdd}
  l_{ij}/l_{ik}\leq 1 /\sin (\pi /6 - \pi /1000)< \sqrt{5}
\end{equation}
 for any two edges in the triangle.
 Without loss of generality, assume that $r_i = 1$. We will prove that $r_j\leq 20$ by contradiction in the following two cases.
	
	If $r_j>20$ and $r_k/r_j\leq 1/5$, then
	$$\frac{l_{ij}^2}{l_{ik}^2} \geq \frac{r_j^2/5 + 2Ir_j+ 1+  4r_j^2/5}{(r_j/5)^2 + 2Ir_j/5+ 1} \geq 5,$$
	which contradicts (\ref{lij bdd}).
	
	If $r_j>20$ and $r_k/r_j> 1/5$, then $r_k>4$ and the inner angle $\theta_{jk}^i$ at $v_i$ is the largest inner angle in $\triangle v_iv_jv_k$.
Note that in this case, we have $I({r_k} + {r_j} - {r_k}{r_j}) + 1<0$, $l_{ij}<\sqrt{I}(r_j+1)$ and $l_{ik}<\sqrt{I}(r_k+1)$.
As a result, by the cosine law, we have
	    $$
	\begin{aligned}
		\cos \theta _{jk}^i &  = \frac{{l_{ij}^2 + l_{ik}^2 - l_{jk}^2}}{{2{l_{ij}}{l_{ik}}}} \\
&= \frac{{I({r_k} + {r_j} - {r_k}{r_j}) + 1}}{{{l_{ij}}{l_{ik}}}} \\
		&  < \frac{{ - I({r_k} + {r_j} + {r_k}{r_j} + 1) + 2I({r_k} + {r_j}) + I + 1}}{{I({r_k} + {r_j} + {r_k}{r_j} + 1)}}
		\\
		& \le   - 1 + \frac{{2I({r_k} + {r_j}) + 2I}}{{I({r_k} + {r_j} + {r_k}{r_j} + 1)}}\\
& \le \frac{{ - 11}}{{21}}.\\
	\end{aligned}
	$$
	This contradicts that the angle bound is $[\pi /6 - \theta_0,  \pi/2 + \theta_0]$.

	
%	$$\cos \theta_{jk}^i = \frac{l_{ij}^2 + l_{ik}^2 - l_{jk}^2}{2l_{ij}l_{ik}} \leq \frac{I( r_k +r_j - r_kr_j ) + 1}{l_{ij}l_{ik}} \leq -1 +\frac{3(r_j+r_k)}{r_jr_k} \leq \frac{-1}{10}. $$
	
	To prove part (b), the definition of $\eta_{ij}^k$ and the formula (\ref{h_ij,k}) implies
	$\eta_{ij}^k = \frac{h_{ij,k}}{l_{ij}} = \frac{r_i^2r_j^2r_kh_k}{Al^2_{ij}},$
	where $A=l_{ij}l_{ik}\sin \theta^i_{jk}$. The sign of $\eta_{ij}^k$ is determined by $h_k$. We will show
	\begin{equation}
	\label{radius}
	    	r_kh_k = (1+I)(1 + I(\frac{r_k}{r_i} +\frac{r_k}{r_j} - 1))\geq \frac{1+I}{4}>0.
	\end{equation}
	We just need to check the case that $r_k/r_i\leq 1$ and $r_k/r_j\leq 1$.  If both $r_k/r_i\geq 1/2$ and $r_k/r_j\geq 1/2$, then $r_kh_k\geq 1+I$. Hence, we only need to consider the situation that  $r_k/r_i<1/2$ or $r_k/r_j<1/2$. By the angle bound and cosine law, we have
	$$-\frac{1}{5(20+I)}l_{ij}^2 \leq l_{jk}^2 + l_{ik}^2 -l_{ij}^2 .$$
	This is equivalent to
		$$I(r_ir_k + r_jr_k - r_ir_j ) \geq - r_k^2 - \frac{1}{10(20+I)}(r_i^2 + r_j^2 + 2Ir_ir_j),$$
which implies
	 $$1 + I(\frac{r_k}{r_i} +\frac{r_k}{r_j} - 1) \geq 1 - \frac{r_k^2}{r_ir_j} - \frac{1}{10(20+I)}(\frac{r_i}{r_j} + \frac{r_j}{r_i} + 2I) \geq \frac{4}{5} - \frac{r_k^2}{r_ir_j},$$
where the results in part (a) of Lemma \ref{ratio} is used in the last inequality.
	Then by the formula (\ref{h_i}) of $h_k$, we have
	$$r_kh_k = (1+I)(1 + I(\frac{r_k}{r_i} +\frac{r_k}{r_j} - 1)) \geq (1+I)(\frac{4}{5} - \frac{r_k^2}{r_ir_j}).$$
	 Therefore, under the assumption that $r_k/r_i\leq 1$ and $r_k/r_j\leq 1$, we have $r_kh_k\geq 3(1+I)/10> (1+I)/4$ when  $r_k/r_i<1/2$ or $r_k/r_j<1/2$.

	The sine law implies that $l_{ij}^2/50\leq A\leq 5l_{ij}^2$. Combining with part (a) of Lemma \ref{ratio}, we can find two constants $M=M(I)$ and $C=C(I)$ such that
	$$M(I) \geq \frac{r_i^2r_j^2 (1+I)(1+40I)}{Al^2_{ij}}\geq \frac{r_i^2r_j^2r_kh_k}{Al^2_{ij}}\geq \frac{r_i^2r_j^2 (1+I)}{4Al^2_{ij}} \geq C(I) > 0.$$ \end{proof}

\end{document}
